\documentclass[10pt,a4paper]{amsart}
\usepackage[margin=19mm]{geometry}
\usepackage{amsmath,amssymb,mathtools}
\usepackage[scr=rsfs,cal=euler]{mathalfa}
\usepackage{xfrac}
\usepackage[unicode,hidelinks,bookmarks=false]{hyperref}
\newcommand{\ie}{i.e.\ }
\newcommand{\cf}{cf.\ }
\newcommand{\resp}{respectively\ }
\newcommand{\Subsection}{\subsection}
\providecommand{\nobreakdash}{\nobreak\hbox{-}}
\setlength{\emergencystretch}{2em}
\multlinegap0pt
\usepackage{xcolor}
\usepackage{tikz}
\usepackage{stackengine}
\usepackage{amssymb}
\newtheorem{cor}{Corollary}
\newtheorem{thm}{Theorem}
\newtheorem{lem}{Lemma}
\title{On Some Properties of Matrices with Entries Defined by Products of $k$-Fibonacci and $k$-Lucas Numbers}
\author{Pedro Fernando Fern\'andez Espinosa, Maritza Liliana Arciniegas Torres, Camilo Andr\'es Acevedo Cadena}
\date{Source: arXiv:2601.12644v2 (CC BY 4.0)}
\begin{document}
\maketitle

\noindent\emph{Source and licence.} On Some Properties of Matrices with Entries Defined by Products of $k$-Fibonacci and $k$-Lucas Numbers. Version arXiv:2601.12644v2 (CC BY 4.0), distributed by arXiv under the Creative Commons Attribution 4.0 International licence, http://creativecommons.org/licenses/by/4.0/, as recorded in the arXiv licence metadata for this exact version and shown on the abstract page https://arxiv.org/abs/2601.12644v2. Full licence notice: \texttt{licenses/arxiv-2601.12644-CC-BY-4.0.txt}.\par\smallskip

\noindent\textit{Editorial scope.} Counted unit: the article's lem charp (source0.tex lines 602--607). The article's complete original proof of it is delivered with it (source0.tex lines 609--802, one \texttt{proof} environment of 194 lines). No mathematics is written, repaired or re-derived: the statement and the proof are the article's own lines, in the article's own notation, and no retained line is substituted or re-flowed.  Retained uncounted as the source-local context the proof needs: lines 305--310; lines 313--320.  Nothing was omitted from any retained span, so the package carries no omission record.  No retained line contains a citation, so the wrapper carries no bibliography and no bibliography entry.  No retained line contains a figure environment, an image-inclusion command or drawing code, so no image is delivered and the package declares no asset dependency.  Editorial formatting only, in addition: an AMS \texttt{amsart} preamble that restates the article's own declarations and macros used by the retained lines, namely the environments \texttt{cor}, \texttt{thm}, \texttt{lem} on separate counters, together with the standard packages those lines need.  The retained environments print in this document as Corollary 1, Theorem 1, Lemma 1 in order of appearance, whereas the article prints the same items with the numbering of its own full document.  The complete native source of the article as distributed in the arXiv e-print for this exact version is bundled unchanged as \texttt{sources/192969/source0.tex}.
\begin{cor}\label{aux1}
For $k=1$, the following identity holds:
\[
F_{2n+1}L_{2n+2}-F_{2n+2}L_{2n+1}=2.
\]
\end{cor}

\begin{thm}\label{aux2}
The following identity holds:
\[
\sum_{i=1}^{n-1}F_{2i}L_{2i-1}+F_{2n-1}L_{2n}
=
2+\sum_{i=1}^{n}F_{2i}L_{2i-1}.
\]
\end{thm}

\begin{lem}\label{charp}
The characteristic polynomial $p(\lambda)$ associated with $A_n$ is given by
\[
p(\lambda)=-(2-\lambda)^{n-1}\left(\lambda-\displaystyle\sum_{i=1}^{n-1}F_{2i}L_{2i-1}+F_{2n-1}L_{2n}\right).
\]
\end{lem}

\begin{proof} We proceed by induction: \\
For $n=2$ we have: 
\begin{align*}
\left|A_n-\lambda I_n\right|&=\left|\left(\begin{array}{ll}
3 & 12 \\
1 & 14
\end{array}\right)-\left(\begin{array}{ll}
\lambda & 0 \\
0 & \lambda
\end{array}\right)\right| \\
&=\left|\left(\begin{array}{cc}
3-\lambda & 12 \\
1 & 14-\lambda
\end{array}\right)\right|\\
& =(3-\lambda)(14-\lambda)-12 =42-3 \lambda-14 \lambda+\lambda^2-12 \\
& =\lambda^2-17 \lambda+30  =(\lambda-2)(\lambda-15) =-(2-\lambda)(\lambda-15)
\end{align*}
On the other hand note that $F_{2}L_{1}+F_{3}L_{4}=1+14=15$. Thus, the base case holds.    

\noindent Now, we assume that the identity holds for some fixed $n\in \mathbb{N}$, we mean that, 

$$
\operatorname{det}\left(A_n-\lambda I_n\right) =-(2-\lambda)^{n-1}\left(\lambda-\sum_{i=1}^{n-1}F_{2 i}L_{2i-1}+F_{2n-1}L_{2 n}\right)
$$
\noindent We want to show that: 
$$
\operatorname{det}\left(A_{n+1}-\lambda I_{n+1}\right)=-(2-\lambda)^n\left(\lambda-\sum_{i=1}^nF_{2i} L_{2 i-1}+F_{2 n+1} L_{2n+2}\right)
$$

\noindent We start by computing $\operatorname{det}\left(B\right)$, where $B=A_{n+1}-\lambda I_{n+1}$, so, we want to compute  

$$\left|\begin{array}{ccccccc}
F_1 L_2-\lambda & F_4 L_3 & F_6 L_5 & F_8 L_7 & \cdots & F_{2 n} L_{2 n-1} & F_{2 n+2} L_{2 n+1} \\
F_2 L_1 & F_3 L_4-\lambda & F_6 L_5 & F_8 L_7 & \cdots & F_{2 n} L_{2 n-1} & F_{2 n+2} L_{2 n+1} \\
F_2 L_1 & F_4 L_3 & F_5L_6-\lambda & F_8 L_7 & \cdots & F_{2 n} L_{2 n-1} & F_{2 n+2} L_{2 n+1} \\
F_2 L_1 & F_4 L_3 & F_6 L_5 & F_7 L_8-\lambda & \cdots & F_{2 n} L_{2 n-1} & F_{2 n+2} L_{2 n+1} \\
\vdots & \vdots & \vdots & \vdots & \ddots & \vdots & \vdots \\
F_2 L_1 & F_4 L_3 & F_6 L_5 & F_8 L_7 & \cdots & F_{2 n-1} L_n-\lambda & F_{2 n+2} L_{2 n+1} \\
F_2 L_1 & F_4 L_3 & F_6 L_5 & F_8 L_7 & \cdots & F_{2 n} L_{2 n-1} & F_{2 n+1} L_{2 n+2}-\lambda
\end{array}\right|$$

\noindent Since $\operatorname{det}\left(A_{n+1}-\lambda I_{n+1}\right)$ can be calculated by expanding by cofactors using any row of $B$ we compute $\operatorname{det}\left(A_{n+1}-\lambda I_{n+1}\right)$ with the last row: 
\begin{align*}
\operatorname{det}(B)
&= \operatorname{det}(A_{n+1}-\lambda I_{n+1}) \\
&= (-1)^{n+2}F_2L_1|M_1|
 + (-1)^{n+3}F_{4}L_{3}|M_2| + \cdots \\
&\quad + (-1)^{2n+1}F_{2n}L_{2n-1}|M_n| \\
&\quad + (-1)^{2n+2}\bigl(F_{2n+1}L_{2n+2}-\lambda\bigr)|M_{n+1}|.
\end{align*}


\noindent where $M_k$ is an $n \times n$ matrix obtained from the matrix $B$ by deleting the $k$-th column and the $(n+1)$-st row, and which can be explicitly described as
$$
M_k=(m_{i,j})=\left\{\begin{array}{lll}
b_{i,j} & \text { if } & j<k \\
b_{i,j+1} & \text {if} & k \leq j \leq n
\end{array}\right.
$$

\noindent when $1\leq k\leq n$ where $b_{i,j}$ denoted the corresponding entry of the matrix $B$ and $1\leq i \leq n$ and $M_{n+1}$ is equal to $A_{n}-\lambda I_n$.

\noindent The above shows that, in order to compute $\operatorname{det}\left(A_{n+1}-\lambda I_{n+1}\right)$, it suffices to determine the value of $M_k$ for $1 \leq k \leq n$. To this end, we consider the matrix $\overline{M}_k$, which is obtained from $M_k$ by applying the admissible row transformations consisting of adding $-1$ times the row $r_k$ to every row $r_t$, for all $t \neq k$. These admissible transformations allow us to describe $\overline{M}_k$ explicitly as follows:

$$
\overline{M_k}=(b_{i,j})= \begin{cases}b_{ij} & \text { si } i=k \\ 
2-\lambda & \text { si } i=j  \text{ and } i<k \\ 
2-\lambda & \text { si } j=i-1 \text { and } i>k \\ 
0 & \text { otherwise }\end{cases}
$$
Thus,
\begin{align*}
\operatorname{det}(B)
&= \operatorname{det}(A_{n+1}-\lambda I_{n+1}) \\
&= (-1)^{n+2}F_2L_1|\overline{M}_1|
 + (-1)^{n+3}F_{4}L_{3}|\overline{M}_2| + \cdots \\
&\quad + (-1)^{2n+1}F_{2n}L_{2n-1}|\overline{M}_n| \\
&\quad + (-1)^{2n+2}\bigl(F_{2n+1}L_{2n+2}-\lambda\bigr)|A_n-\lambda I_n|.
\end{align*}

\noindent To compute $|\overline{M}_k|$ we again expand by cofactors using the row $k$ of the corresponding $\overline{M}_k$, that is, $|\overline{M}_k|=\displaystyle \sum_{j=1}^{n}b_{kj}|P_j|$ where $P_j$ is a matrix of size $(n-1)\times(n-1)$ with the following form
\[
P_j = (P_{s,t})=
\begin{cases}
2 - \lambda, & \text{if } s = t \text{ and } s < j, \\
2 - \lambda, & \text{if } t = s - 1 \text{ and } s > j, \\
0,           & \text{otherwise}.
\end{cases}
\]

\noindent Note that the column $n-1$ is always $0$ for all $P_j$ with $1\leq j \leq n-1$. By properties of determinants $|P_j|=0$ for $1\leq j \leq n-1$. Thus, $|\overline{M}_k|=(-1)^{n+k}b_{k,n}|P_n|$, but $P_n=I_{n-1}(2-\lambda)$ which implies that: 
\[|\overline{M_{k}}|=(-1)^{n+k}F_{2n+2}L_{2n+1}(2-\lambda)^{n-1}\]
So, applying the induction hypothesis we get: 

\begin{align}
\operatorname{det}(B)
&= \operatorname{det}(A_{n+1}-\lambda I_{n+1}) \notag \\
&= (-1)^{n+2}F_{2}L_{1}\cdot(-1)^{n+1}F_{2n+2}L_{2n+1}(2-\lambda)^{n-1} \notag \\
&\quad + (-1)^{n+3}F_{4}L_{3}\cdot(-1)^{n+2}F_{2n+2}L_{2n+1}(2-\lambda)^{n-1} \notag \\
&\quad + \cdots \notag \\
&\quad + (-1)^{2n+1}F_{2n}L_{2n-1}\cdot(-1)^{2n}F_{2n+2}L_{2n+1}(2-\lambda)^{n-1} \notag \\
&\quad + (-1)^{2n+2}\bigl(F_{2n+1}L_{2n+2}-\lambda\bigr) \cdot -(2-\lambda)^{n-1}
\left(\lambda-\sum_{i=1}^{n-1}F_{2i}L_{2i-1}+F_{2n-1}L_{2n}\right)\notag
\end{align}

\noindent which can be simplify and factorize by 
\[
-(2-\lambda)^{n-1}\Bigl(
F_{2n+2}L_{2n+1}\Bigl( \sum_{i=1}^{n} F_{2i}L_{2i-1} \Bigr)
+\bigl(F_{2n+1}L_{2n+2}-\lambda\bigr)
\bigl(\lambda-\sum_{i=1}^{n-1}F_{2i}L_{2i-1}+F_{2n-1}L_{2n}\bigr)
\Bigr).
\]

\noindent Simplifying and reordering 

\[
\begin{aligned}
-(2-\lambda)^{n-1}\Bigl(
&-\lambda^{2}
+ \lambda\Bigl(
F_{2n+1}L_{2n+2}
+ \sum_{i=1}^{n-1} F_{2i}L_{2i-1}
+ F_{2n-1}L_{2n}
\Bigr) \\
&- F_{2n+1}L_{2n+2}
\Bigl(
\sum_{i=1}^{n-1} F_{2i}L_{2i-1}
+ F_{2n-1}L_{2n}
\Bigr) \\
&+ F_{2n+2}L_{2n+1}
\Bigl(
\sum_{i=1}^{n} F_{2i}L_{2i-1}
\Bigr)
\Bigr).
\end{aligned}
\]




\noindent Now, we use Corollary \ref{aux1} and Theorem \ref{aux2} to simplify the last expression: 

\begin{align*}
\operatorname{det}(B)&=-(2-\lambda)^{n-1} \Bigl( -\lambda^2+\lambda\left(F_{2 n+1} L_{2 n+2}+\sum_{i=1}^{n-1} F_{2 i} L_{2 i-1}+F_{2 n-1} L_{2 n}\right)\\
&-F_{2 n+1} L_{2 n+2} \cdot\left(\sum_{i=1}^{n-1} F_{2 i} L_{2 i-1}+F_{2 n-1} \cdot L_{2 n}\right)\\ 
&+F_{2 n+2} L_{2 n+1}\left(\sum_{i=1}^n F_{2 i} L_{2 i-1}\right)\Bigr)\\
&= -(2-\lambda)^{n-1} \Bigl(-\lambda^2+\lambda\left(F_{2 n+1} L_{2 n+2}+\sum_{i=1}^{n-1} F_{2i} L_{2 i-1}+F_{2 n-1} L_{2 n}\right)\\
&-\Bigl(\left(2+F_{2 n+2} L_{2 n+1}\right) \cdot\left(2+\sum_{i=1}^n F_{2 i} L_{2 i-1}\right)\Bigr) \\
&  +F_{2 n+2} L_{2 n+1} \cdot\left(\sum_{i=1}^n F_{2 i} L_{2 i-1}\right)\Bigr)
\end{align*}
\normalcolor
The second factor is equivalent to 

\begin{align*}
&=-\lambda^2+\lambda(F_{2n+1}L_{2n+2}+\sum_{i=1}^{n-1}F_{2i}L_{2i-1}+F_{2n-1}L_{2n})\\
&-\Bigl(2\cdot2+2\sum_{i=1}^{n}F_{2i}L_{2i-1}+2F_{2n+2}L_{2n+1}+F_{2n+2}L_{2n+1}\cdot\sum_{i=1}^{n}F_{2i}L_{2i-1}\Bigr)\\
&+F_{2n+2}L_{2n+1}\Bigl(\sum_{i=1}^{n}F_{2i}L_{2i-1}\Bigr)
\end{align*}
%%%%%%%%%%%%%%%%%%%%%%%%
\begin{align*}
&=-\lambda^2+\lambda(F_{2n+1}L_{2n+2}+\sum_{i=1}^{n-1}F_{2i}L_{2i-1}+F_{2n-1}L_{2n})\\
&-2\cdot2-2\sum_{i=1}^{n}F_{2i}L_{2i-1}-2F_{2n+2}L_{2n+1}-F_{2n+2}L_{2n+1}\cdot\sum_{i=1}^{n}F_{2i}L_{2i-1}\\
&+F_{2n+2}L_{2n+1}\Bigl(\sum_{i=1}^{n}F_{2i}L_{2i-1}\Bigr)
\end{align*}
%%%%%%%%%%%%%%%%%%%%%%%
\begin{align*}
&=-\lambda^2+\lambda(F_{2n+1}L_{2n+2}+\sum_{i=1}^{n-1}F_{2i}L_{2i-1}+F_{2n-1}L_{2n})\\
&-2\Bigl( 2+\sum_{i=1}^{n}F_{2i}L_{2i-1}+F_{2n+2}L_{2n+1}\Bigr)
\end{align*}
%%%%%%%%%%%%%%%%%%%%%%%
\begin{align*}
&=-\lambda^2+\lambda(F_{2n+1}L_{2n+2}+\sum_{i=1}^{n-1}F_{2i}L_{2i-1}+F_{2n-1}L_{2n})\\
&-2\Bigl( \sum_{i=1}^{n}F_{2i}L_{2i-1}+F_{2n+1}L_{2n+2}\Bigr)
\end{align*}
\noindent Note that by Theorem \ref{aux2} 

\[\sum_{i=1}^{n-1}F_{2i}L_{2i-1}+F_{2n-1}L_{2n}+F_{2n+1}L_{2n+2}=2+\sum_{i=1}^{n}F_{2i}L_{2i-1}+F_{2n+1}L_{2n+2}\]

\noindent So the expression becomes

\[-\lambda^2+\lambda(2+\sum_{i=1}^{n}F_{2i}L_{2i-1}+F_{2n+1}L_{2n+2})\\
-2\Bigl( \sum_{i=1}^{n}F_{2i}L_{2i-1}+F_{2n+1}L_{2n+2}\Bigr)\]
\noindent It can be factor as

\[(2-\lambda)\Bigl( \lambda - \sum_{i=1}^{n} F_{2i}L_{2i-1}+F_{2n+1}L_{2n+2} \Bigr)\]

\noindent Thus 
\[\operatorname{det}(A_{n+1}-\lambda I_{n+1})=-(2-\lambda)^{n-1}\cdot(2-\lambda)\cdot\Bigl( \lambda - \sum_{i=1}^{n} F_{2i}L_{2i-1}+F_{2n+1}L_{2n+2} \Bigr)\]
So, 
\[\operatorname{det}(A_{n+1}-\lambda I_{n+1})=-(2-\lambda)^{n}\cdot\Bigl( \lambda - \sum_{i=1}^{n} F_{2i}L_{2i-1}+F_{2n+1}L_{2n+2} \Bigr)\]

\noindent This completes the inductive step, and thus the proof.
\end{proof}

\end{document}
